\section{A general scheme of agentic AutoML}
\label{sec:scheme}

% Added 26 October 2026 after Boukhanovsky's comment 7: the general part
% before the harnesses.  The published systems in Table 1 were checked
% against their papers (arXiv versions: AutoKaggle v3, DS-Agent v5,
% R&D-Agent v2, KompeteAI v3).

This section describes an agentic AutoML system in general terms. It
defines the outcomes of a trial, names the parts of a system and the
role of each, and places published systems in one design space. The three
hypotheses are then stated as claims about these parts.

\subsection{A trial and its outcomes}
\label{sec:outcomes}

A task consists of training data, a description, a metric and a time
budget. A trial is one run of a system on one task. It ends with a
submission, a set of predictions for the test data, or with nothing when the
system stops or runs out of time. Three outcomes describe a trial.
\emph{Reliability} is the share of trials that end with a submission.
\emph{Accuracy} is the test score of a submission and is measured only
where one exists. \emph{Cost} is the wall-clock time and the tokens an
trial consumes. A system can gain in one outcome and lose in another, so
the paper reports the three separately.

\subsection{The parts of a system}
\label{sec:parts}

A data scientist works in a loop: explore the data, plan, write code, run
it, read the result and repeat until a submission is ready. An agentic
AutoML system splits this loop among five parts
(Figure~\ref{fig:scheme}).

The \emph{language model} chooses the next action from what it has seen so
far. It is needed because the space of solutions is open: any code the
model can write is a candidate, and no grid of pipelines is fixed in
advance.

The \emph{harness} turns the choices of the model into work. Each of its
four components answers one limit of a language model.
\begin{itemize}
\setlength{\itemsep}{1pt}
\item \emph{Actions} let the text of the model run: a shell, code
  execution, calls to a library.
\item \emph{Control} sets who plans. Designers fix the stages in advance,
  an agent writes a plan at run time, or a single agent chooses each step
  as it goes. A task of an hour takes many model calls, and control orders
  them.
\item The \emph{execution loop} returns the output and the errors of the
  code to the model. Generated code often fails on the first run, and the
  loop decides how many repairs follow and who makes them.
\item \emph{Memory} sets what passes from one step to the next: one growing
  context, written reports between agents, or fixed input and output
  contracts. The context of a model is finite, and memory decides what
  stays in it.
\end{itemize}

\emph{Knowledge about the procedure} tells the model which tools to use and
how to work. It comes from the parametric memory of the model, from
retrieved earlier solutions~\citep{guo2024dsagent,kulibaba2025kompeteai},
from library documentation built into the system~\citep{fang2025mlzero} or
from an instruction file written by experts.

The \emph{model builder} fits the final model. The agent can write the
model in its own code, or call an AutoML library as a tool and let it
search and tune. This part joins agentic AutoML with the AutoML frameworks
of Section~\ref{sec:intro}.

The \emph{environment} is the container with its libraries, the hardware
and the time limit of a trial. It sets which tools exist and how long a
fit may run.

\subsection{Published systems in the scheme}
\label{sec:space}

Table~\ref{tab:designpoints} places eight published systems and the three
harnesses of this study in the scheme. The published systems are described
from their papers. Their scores are not compared with ours: they were
obtained with other language models, such as GPT-4o in AutoKaggle and
AutoML-Agent, Claude~3.7 Sonnet in MLZero and GPT-5 and o3 in R\&D-Agent,
and on other benchmarks.

The table shows that the parts vary independently in practice. AutoKaggle
and DS-Agent both fix the stages. The first fits models through its own
tool library, and the second writes the model in its own code. Six of the
eight systems write the model themselves, and MLZero and LightAutoDS-Tab
can call an AutoML library. Four systems retrieve external knowledge, and
LightAutoDS-Tab and AutoKaggle also use text written by people. Two systems
that share one pattern of control can therefore differ in every other part,
and a comparison of complete systems changes several parts at once.

\begin{table*}[t]
\centering
\caption{Published agentic AutoML systems and the three harnesses of this
study on the parts of the scheme. Published systems are described from
their papers.}
\label{tab:designpoints}
\footnotesize
\setlength{\tabcolsep}{3pt}
\renewcommand{\arraystretch}{1.1}
\newcolumntype{L}[1]{>{\raggedright\arraybackslash}p{#1}}
\begin{tabular}{@{}L{2.3cm}L{2.6cm}L{0.9cm}L{2.2cm}L{2.3cm}L{2.4cm}L{2.3cm}L{1.6cm}@{}}
\toprule
System & Control & Agents & Memory & Execution loop & Model builder & Knowledge & Inputs \\
\midrule
\rowcolor{stripe}AutoKaggle~\citep{li2024autokaggle} & fixed phases, LLM plan within each & 5 & reports & bounded debugging, unit tests & tool library on scikit-learn, own code & tool documentation, rules by people & tabular \\
DS-Agent~\citep{guo2024dsagent} & fixed loop, LLM plan per step & roles & experiment log & bounded debugging & own code & Kaggle case bank & text, time series, tabular \\
\rowcolor{stripe}AutoML-Agent~\citep{trirat2025automlagent} & plans generated by an LLM & 5 & JSON contract, code skeleton & verification and revision & own code & web, arXiv, dataset search & five modalities \\
MLZero~\citep{fang2025mlzero} & fixed modules & 9 & context, error summaries & bounded retries & AutoGluon or own code & library tutorials & several \\
\rowcolor{stripe}LightAutoDS-Tab~\citep{lapin2025lightautods} & router, LLM plan on a code skeleton & 7 & messages, code template & improvement loop & own code, LightAutoML or FEDOT & prompts by designers & tabular \\
R\&D-Agent~\citep{yang2025rdagent} & search over candidates & 2 & exploration graph & many candidates, selection & own code & none & several \\
\rowcolor{stripe}KompeteAI~\citep{kulibaba2025kompeteai} & tree search over fixed stages & 8 & solution tree & many candidates, pruning & own code & Kaggle and arXiv retrieval & not stated \\
AIRA~\citep{toledo2025aira} & search policy & 1 & search graph & many candidates, selection & own code & none & not stated \\
\midrule
\multicolumn{8}{@{}l}{\emph{The three harnesses of this study}} \\
\rowcolor{stripe}FEDOT.LLM & fixed stages & 1 & fixed contracts & bounded retries & FEDOT, fixed & library built in & tabular \\
AutoDS-Tools & plan generated by agents & 6 & written reports & per-agent debugging & own code, or LightAutoML when the file names it & instruction file, on or off & any \\
\rowcolor{stripe}Terminus-2 & step by step & 1 & one growing context & open loop & own code & none, file appended in an ablation & any \\
\bottomrule
\end{tabular}
\end{table*}

\subsection{The hypotheses in terms of the scheme}
\label{sec:hyp-scheme}

\ref{hyp:harness} concerns the harness. With the language model,
the knowledge and the model builder fixed, control, the execution loop and
memory decide how a trial proceeds. They should change reliability and
cost. The final model comes from the same language model, so accuracy
should change little. The hypothesis is stated for harnesses that do not
search over candidate solutions, for the reason given in
Section~\ref{sec:why-three}.

\ref{hyp:builder} concerns the model builder. The agent's own
code and an AutoML library differ in whether a search over models and
hyperparameters is run, so the model builder should decide accuracy. A
harness that fixes a library and an instruction that names one both act
through this part.

\ref{hyp:interaction} concerns knowledge and the language model.
An instruction assumes a control flow and an execution loop, and a language
model fills the time limit at its own pace. Their effect on reliability
should therefore depend on the harness and the time limit, and can change
sign from one harness to another.

Each hypothesis is a claim about parts of the scheme and outcomes of an
trial, so it can be tested on any system placed in it.

\subsection{Why three harnesses}
\label{sec:why-three}

The test needs harnesses that differ on the parts of the scheme while the
language model, the knowledge and the environment stay fixed. The study
therefore compares parts of systems, and each pattern is represented by one
harness. The three harnesses run in one framework,
Harbor~\citep{harbor2026}, with the same language model, and AutoDS-Tools
has an instruction file that can be switched part by part. Running the
published systems as their papers describe them would change the language
model, the knowledge and the protocol together, the confound described in
Section~\ref{sec:intro}.

On control the three harnesses take three values of
Table~\ref{tab:designpoints}: fixed stages in FEDOT.LLM, a plan generated
by agents in AutoDS-Tools and step-by-step choice in Terminus-2. These
follow the division of LLM systems into workflows, multi-agent systems and
autonomous agents~\citep{anthropic2024agents}. On memory they take fixed
contracts, written reports and one growing context. On the model builder
they take a fixed library, a library named by the instruction file and the
agent's own code.

Five of the eight published systems fix the stages or let agents generate a
plan, as FEDOT.LLM and AutoDS-Tools do. The other three search over many
candidate solutions and keep the best on
validation~\citep{yang2025rdagent,kulibaba2025kompeteai,toledo2025aira}.
Such a harness runs the search that \ref{hyp:builder} attributes
to the model builder, so \ref{hyp:harness} is stated and tested
for harnesses without it. On the model builder the published systems take
the two values our harnesses take: their own code, or an AutoML library in
MLZero and LightAutoDS-Tab. Retrieval of external knowledge, used by four
systems, lies outside the study.

FEDOT.LLM fixes the model builder to the FEDOT library. Without the
instruction file, \ref{hyp:harness} is therefore tested cleanly
on AutoDS-Tools and Terminus-2, which both write their own code. FEDOT.LLM
enters the same comparison with a different model builder. FEDOT runs its
search only when fitting the default pipeline leaves time, so
\ref{hyp:builder} predicts that FEDOT.LLM gains on the tasks
where the search runs.
